\documentclass[11pt,a4paper]{article}

\usepackage[margin=1in]{geometry}
\usepackage[T1]{fontenc}
\usepackage[utf8]{inputenc}
\usepackage{lmodern}
\usepackage{xcolor}
\usepackage{hyperref}
\usepackage{enumitem}
\usepackage{fancyvrb}
\usepackage{titlesec}
\usepackage{longtable}

\hypersetup{
  colorlinks=true,
  linkcolor=blue,
  urlcolor=blue,
  pdftitle={Git Feature Workflow Course},
  pdfauthor={ManhwaRank: The Archive}
}

\titleformat{\section}{\large\bfseries}{\thesection}{0.6em}{}
\titleformat{\subsection}{\normalsize\bfseries}{\thesubsection}{0.6em}{}

\title{Git Feature Workflow Course \\ \large Practical Git for ManhwaRank: The Archive}
\author{Project Training Notes}
\date{\today}

\begin{document}
\maketitle
\tableofcontents
\newpage

\section{Course Goal}
This course teaches a complete, repeatable Git workflow for adding new features safely to your project.

After this course, you should be able to:
\begin{itemize}[leftmargin=1.5em]
  \item create feature branches correctly,
  \item commit cleanly with meaningful history,
  \item keep your branch updated with \texttt{main},
  \item open and merge pull requests,
  \item resolve merge conflicts,
  \item revert mistakes without panic,
  \item deploy confidently after merge.
\end{itemize}

\section{Mental Model}
Think of Git as three layers:
\begin{enumerate}[leftmargin=1.5em]
  \item \textbf{Working tree}: files you edit.
  \item \textbf{Staging area} (index): files selected for next commit.
  \item \textbf{Repository history}: commits already saved.
\end{enumerate}

\subsection{Golden Rules}
\begin{itemize}[leftmargin=1.5em]
  \item Never build features directly on \texttt{main}.
  \item One feature or fix per branch.
  \item Small commits with clear messages are easier to review and debug.
  \item Always pull latest \texttt{main} before branching and before merging.
  \item Do not commit secrets (\texttt{.env}, tokens, API keys).
\end{itemize}

\section{One-Time Setup}
Run these once on your machine.

\subsection{Identity}
\begin{Verbatim}[fontsize=\small]
git config --global user.name "Your Name"
git config --global user.email "you@example.com"
\end{Verbatim}

\subsection{Helpful Defaults}
\begin{Verbatim}[fontsize=\small]
git config --global init.defaultBranch main
git config --global pull.rebase false
git config --global fetch.prune true
git config --global core.autocrlf true
\end{Verbatim}

\subsection{Verify Remote}
\begin{Verbatim}[fontsize=\small]
git remote -v
\end{Verbatim}

Expected remote style:
\begin{Verbatim}[fontsize=\small]
origin  https://github.com/<your-org-or-user>/<repo>.git (fetch)
origin  https://github.com/<your-org-or-user>/<repo>.git (push)
\end{Verbatim}

\section{Project Workflow Standard}
Use this exact cycle for each new feature.

\subsection{Step 1: Start Clean on Main}
\begin{Verbatim}[fontsize=\small]
cd C:\Users\Adill\Documents\project_manhwa_v2\manhwa-aggregator
git checkout main
git pull origin main
git status
\end{Verbatim}

You want: \texttt{working tree clean}.

\subsection{Step 2: Create a Feature Branch}
Branch naming pattern:
\begin{itemize}[leftmargin=1.5em]
  \item \texttt{feature/<short-name>}
  \item \texttt{fix/<short-name>}
  \item \texttt{chore/<short-name>}
\end{itemize}

Example:
\begin{Verbatim}[fontsize=\small]
git checkout -b feature/cors-env-origins
\end{Verbatim}

\subsection{Step 3: Build in Small Increments}
As you code, check changes often:
\begin{Verbatim}[fontsize=\small]
git status
git diff
\end{Verbatim}

Stage only relevant files:
\begin{Verbatim}[fontsize=\small]
git add backend/main.py
git add DEPLOY_FREE.md
\end{Verbatim}

Commit with a strong message:
\begin{Verbatim}[fontsize=\small]
git commit -m "Add env-based CORS origin configuration"
\end{Verbatim}

Good commit message pattern:
\begin{itemize}[leftmargin=1.5em]
  \item Imperative verb first: Add, Fix, Refactor, Update, Remove.
  \item Say what changed, not how hard it was.
  \item Keep first line around 50--72 chars if possible.
\end{itemize}

\subsection{Step 4: Push Branch}
\begin{Verbatim}[fontsize=\small]
git push -u origin feature/cors-env-origins
\end{Verbatim}

\subsection{Step 5: Open Pull Request}
On GitHub:
\begin{enumerate}[leftmargin=1.5em]
  \item Open PR from \texttt{feature/...} into \texttt{main}.
  \item Add clear title and summary.
  \item Verify CI checks pass.
  \item Merge when green.
\end{enumerate}

\subsection{Step 6: Sync Local After Merge}
\begin{Verbatim}[fontsize=\small]
git checkout main
git pull origin main
git branch -d feature/cors-env-origins
git push origin --delete feature/cors-env-origins
\end{Verbatim}

\section{Daily Command Cheatsheet}
\begin{longtable}{p{0.33\textwidth} p{0.6\textwidth}}
\textbf{Command} & \textbf{What it does} \\
\hline
\texttt{git status} & Show changed, staged, and untracked files. \\
\texttt{git diff} & Show unstaged content changes. \\
\texttt{git diff --staged} & Show staged changes for next commit. \\
\texttt{git add <file>} & Stage a specific file. \\
\texttt{git add -p} & Stage selected hunks interactively. \\
\texttt{git commit -m "msg"} & Create a commit from staged files. \\
\texttt{git log --oneline -n 10} & Show recent commit history. \\
\texttt{git pull origin main} & Update local main with remote main. \\
\texttt{git checkout -b <branch>} & Create and switch to a new branch. \\
\texttt{git push -u origin <branch>} & Push branch and set upstream. \\
\end{longtable}

\section{Feature Development Example (Realistic)}
Goal: add environment-based CORS allow list.

\subsection{Branch}
\begin{Verbatim}[fontsize=\small]
git checkout main
git pull origin main
git checkout -b feature/cors-env-origins
\end{Verbatim}

\subsection{Implement + Validate}
\begin{Verbatim}[fontsize=\small]
uvicorn backend.main:app --reload --port 8000
\end{Verbatim}

If startup is successful, stop with Ctrl+C.

\subsection{Commit}
\begin{Verbatim}[fontsize=\small]
git add backend/main.py DEPLOY_FREE.md
git commit -m "Support ALLOWED_ORIGINS env var for CORS"
git push -u origin feature/cors-env-origins
\end{Verbatim}

\section{How to Handle Mistakes}
\subsection{Case A: Wrong file staged}
\begin{Verbatim}[fontsize=\small]
git restore --staged <file>
\end{Verbatim}

\subsection{Case B: Edited file but want to discard local changes}
\begin{Verbatim}[fontsize=\small]
git restore <file>
\end{Verbatim}

\subsection{Case C: Last commit message is wrong}
\begin{Verbatim}[fontsize=\small]
git commit --amend -m "Better commit message"
\end{Verbatim}

\subsection{Case D: Need to undo one bad commit safely}
\begin{Verbatim}[fontsize=\small]
git revert <commit-hash>
\end{Verbatim}

\textbf{Use \texttt{revert} for shared branches.} It creates a new commit that undoes changes without rewriting history.

\section{Merge Conflict Playbook}
When \texttt{main} changed the same lines as your branch:

\subsection{Update Branch}
\begin{Verbatim}[fontsize=\small]
git checkout feature/your-branch
git fetch origin
git merge origin/main
\end{Verbatim}

\subsection{Resolve}
\begin{enumerate}[leftmargin=1.5em]
  \item Open conflicted files.
  \item Remove conflict markers:\\
  \texttt{<<<<<<<}, \texttt{=======}, \texttt{>>>>>>>}.
  \item Keep the correct final code.
  \item Run app/tests.
\end{enumerate}

\subsection{Finalize}
\begin{Verbatim}[fontsize=\small]
git add <resolved-file-1> <resolved-file-2>
git commit
git push
\end{Verbatim}

\section{Git and Deployment (Your Current Stack)}
\subsection{Triggering Deploys}
\begin{itemize}[leftmargin=1.5em]
  \item Push to \texttt{main} triggers Render and Vercel auto deploys.
  \item If build cache is problematic, use \texttt{Redeploy without cache}.
\end{itemize}

\subsection{Important for This Repo}
\begin{itemize}[leftmargin=1.5em]
  \item Keep \texttt{node\_modules} out of Git.
  \item Keep \texttt{.env} out of Git.
  \item Keep frontend root directory correctly configured in Vercel.
  \item Keep backend root directory correctly configured in Render.
\end{itemize}

\section{Recommended Team Policy (Simple but Strong)}
\begin{enumerate}[leftmargin=1.5em]
  \item No direct coding on \texttt{main}.
  \item Every change goes through feature branch + PR.
  \item CI must pass before merge.
  \item One logical feature per PR.
  \item Use descriptive commit messages.
\end{enumerate}

\section{Weekly Hygiene}
Once a week:
\begin{Verbatim}[fontsize=\small]
git checkout main
git pull origin main
git branch --merged
\end{Verbatim}

Delete merged local branches you no longer need.

\section{Troubleshooting FAQ}
\subsection*{Q1: \texttt{remote origin already exists}}
\begin{Verbatim}[fontsize=\small]
git remote set-url origin https://github.com/<owner>/<repo>.git
\end{Verbatim}

\subsection*{Q2: 403 permission denied on push}
Usually wrong cached GitHub account.
\begin{Verbatim}[fontsize=\small]
@"
protocol=https
host=github.com
"@ | git credential-manager erase
\end{Verbatim}
Then push again and sign in with the correct account.

\subsection*{Q3: Vercel build exits with 126}
Likely executable permission issue from tracked \texttt{node\_modules}. Remove tracked dependency folders and redeploy without cache.

\section{30-Second Workflow Summary}
For every new feature:
\begin{Verbatim}[fontsize=\small]
git checkout main
git pull origin main
git checkout -b feature/<name>
# edit files
# run app/tests
git add <relevant-files>
git commit -m "Add <feature>"
git push -u origin feature/<name>
# open PR, merge
git checkout main
git pull origin main
\end{Verbatim}

\section{Final Advice}
Consistency beats cleverness. If you always follow the same branch-commit-PR cycle, Git becomes predictable, safe, and fast.

Your next practical step:
\begin{enumerate}[leftmargin=1.5em]
  \item Pick one small feature.
  \item Do it using this exact workflow.
  \item Repeat until it becomes muscle memory.
\end{enumerate}

\end{document}
